%ist Arsch, ist aber auch noch oke am Anfang.

\chapter{Einleitung}\label{einleitung}
\section{Motivation}
Die Vorhersage von Aktienkursen stellt eine zentrale Herausforderung der Finanzökonomie dar. Trotz umfangreicher Forschung bleibt die präzise Prognose zukünftiger Kursbewegungen aufgrund der Komplexität und Volatilität der Märkte schwierig. In den letzten Jahren wurden vermehrt mathematische und datengetriebene Methoden untersucht, um Muster in Finanzkursen zu identifizieren und Vorhersagen zu verbessern. Da ein zuverlässiger und konsistenter Algorithmus zur Vorhersage von Aktienkursen potenziell erhebliche finanzielle Vorteile bieten würde, ist das Interesse an der Entwicklung solcher Verfahren entsprechend groß.
\section{Probleme}
Ein zentrales theoretisches Problem bei der Vorhersage von Aktienkursen ergibt sich aus der \ac{EMH}. Die \ac{EMH} besagt, dass alle verfügbaren Informationen bereits vollständig in den aktuellen Marktpreisen enthalten sind, sodass zukünftige Kursbewegungen per Definition nicht systematisch vorhergesagt werden können. Unter dieser Annahme verhalten sich Preisänderungen wie ein zufälliger Prozess, der durch keinen Algorithmus vorhergesagt werden kann und die Machbarkeit des Ziels dieser Arbeit in Frage stellt \cite{malkiel1989efficient}.
Trotz der theoretischen Einschränkungen durch die Markteffizienzhypothese existieren empirische Beispiele erfolgreicher quantitativer Handelsstrategien. Renaissance Technologies nutzt seit Jahrzehnten fast ausschließlich mathematische Modelle und Algorithmen zur Kursprognose und erzielt mit dem Medallion Fund Renditen in Höhe von 66\% p.a. Dies zeigt, dass systematische Modelle wiederkehrende Muster in Finanzdaten sehr wohl erkennen und wirtschaftlich nutzen können. Hierbei sei anzumerken, dass die genaue Funktionsweise der verwendeten Algorithmen selbstverständlich unbekannt ist. Quantitative Fonds agieren als "Black-Box", was die wissenschaftliche Nachvollziehbarkeit erschwert und somit kein gesichertes Gegenbeispiel zur \ac{EMH} ist \cite{burton2016inside}. Des Weiteren gibt es empirische, öffentliche Analysen, die zeigen, dass nicht-lineare Vorhersagemethoden durchaus Gewinne erzielen können \cite{Sewell2012_EMH}.
Neben der Frage nach der allgemeinen Machbarkeit ergeben sich eine Vielzahl praktischer Probleme, die die Entwicklung erschweren. Zunächst enthalten Kursdaten viel Rauschen. Es ist schwer den tatsächlichen Signalanteil des Preises herauszufiltern, um zu den Trend zu erkennen. Des Weiteren ist es eine große Schwierigkeit, den korrekten Komplexitätsgrad des Modells zu finden. Zu komplexe Modelle \glqq overfitten\grqq, passen sich also zu stark an historische Daten an und generalisieren somit schlecht auf neuen, realen Daten. Dies betrifft vor allem Deep Learning Modelle, welche zu viele Freiheitsgrade haben, und somit den vergangenenen Kurs oder das Rauschen selbst abbilden und den Trend dabei nicht erkennen. Auf der anderen Seite sind zu simple Modelle nicht in der Lage nicht-lineare Zusammenhänge zu erkennen und generalisieren zu stark \cite{PrakashTuoDing2023_temporalOverfitting}.

\section{Ziel der Arbeit}

Das Ziel dieser Arbeit ist das Entwickeln einer Strategie zur Auswahl von rentablen Aktien aus einem Aktienuniversum und das Entwickel und Testen von Vorhersagemodellen auf diesen gewählten Aktien. Anschließend wird eine Marktstrategie entwickelt und die Modelle werden in einem realistischen Backtest evaluiert. Es soll eine Testumgebung geschaffen werden, mit welcher verschiedene Analysemodelle getestet und auf realen, vergangenen Daten ausprobiert werden können. 
